\section{Proof of integral Donovan finiteness}\label{sec:assembly}

Fix $p$ and $M$.  The previous sections supply three inputs for the
last step: finitely many integral basic identity orders, finite based
restrictions on their $p$-subgroups, and the integral kernel-removal
theorem.  We check these inputs for an arbitrary finite-group block.

\begin{proof}[Proof of Theorem~\ref{thm:main}]
Let $B$ be a block of $\cO G$ with defect order at most $M$.
\emph{Reduction and the identity order.}
The integral An--Eaton reduction gives a reduced pair
$(H,B_H)$ with $B_H$ $\cO$-Morita equivalent to $B$ and with
the same defect group.  Apply
Proposition~\ref{prop:structure}.  Its integral dual recovery
realizes $B_H$, up to full corners, as a block summand of a
crossed order on an identity block $B_0$, with grading group
$Q$ satisfying a uniform bound on $[Q:O_{p'}(Q)]$.

Compress by a basic idempotent of $B_0$.  As in the proof of
Theorem~\ref{thm:based-sylow}, correction of the homogeneous
units makes the resulting full corner a crossed order on
one of the finite list $R_1,\ldots,R_t$.  Its outer action
lies in the finite group $\Out_{\cO}(R_i)$, by
Lemma~\ref{lem:HH} or Proposition~\ref{prop:structure}(iv).

\emph{The restrictions and the kernel.}
Theorem~\ref{thm:based-sylow} supplies finitely many based
restrictions on every possible $p$-subgroup.  Its proof used the exact
integral comparison on the actual extension blocks, followed by
fixed-total-order rigidity; hence these are the based restrictions
required in the kernel theorem.

All three hypotheses of Theorem~\ref{thm:kernel} now hold for this
family over $R_i$: bounded $[Q:O_{p'}(Q)]$, finite outer-action
image, and finite based $p$-subgroup restrictions.  Their block
summands therefore have finitely many integral Morita classes.

\emph{The finite union.}
Taking
the finite union over $i$ gives a list independent of $G$
and $b$.  The full-corner equivalences and the preliminary
integral reduction place the original $B$ in that list.
This proves the theorem.
\end{proof}

\begin{corollary}\label{cor:reduction}
For every $p,M$, there is a finite list of integral basic
orders $R_1',\ldots,R_v'$ such that the basic algebra of
every block of $kG$ of defect order at most $M$ is
isomorphic to $k\otimes_{\cO}R_i'$ for some $i$.
In particular the blocks over $k$ of bounded defect have
finitely many $k$-linear Morita classes.
\end{corollary}

\begin{proof}
Every block idempotent over $k$ has its unique central
idempotent lift to $\cO G$.  Lift a basic idempotent in
that block.  Its corner is an integral basic order, and
reduction gives the original basic algebra.  Theorem~\ref{thm:main}
makes the integral basic orders a finite list, proving
both assertions.
\end{proof}

The residue-field finiteness in Corollary~\ref{cor:reduction} was
already established by the field companion and supplied the Cartan
input in Proposition~\ref{prop:criteria}.  The corollary now identifies
a finite list of integral basic orders whose reductions represent
those field algebras.  The proof of that stronger conclusion has used
the explicit integral constructions throughout; it has not required
lifting an arbitrary field Morita equivalence.

\subsection{Scalar extension to a fixed complete coefficient ring}

\begin{proof}[Proof of Corollary~\ref{cor:complete-dvr}]
Keep $k=\overline{\F}_p$ and $\cO=W(k)$, and fix $\mathcal R$ and
its residue field $\ell$ as in the corollary.  Choose an embedding
$\iota:k\hookrightarrow\ell$.  Since $\ell$ is perfect, $W(\ell)$ is
a Cohen ring.  The coefficient-ring theorem
\cite[Theorem~10.160.8, Tag~032A]{StacksCohen} gives a local map
$W(\ell)\to\mathcal R$ inducing the identity on $\ell$.  It is
injective: any nonzero ideal of the DVR $W(\ell)$ contains a power of
$p$, whereas $\mathcal R$ has characteristic zero.  Composing with the
Witt-vector map $W(\iota)$, which is injective on Witt coordinates,
gives a fixed local embedding
\[
 \cO=W(k)\xrightarrow{W(\iota)}W(\ell)\hookrightarrow\mathcal R
\]
whose residue map is $\iota$.  Only this coefficient embedding is
used; no finiteness of $\mathcal R$ over $\cO$ is needed.

We recall the residue-field argument of
\cite[Section~2.1]{OpenAIAlperinWeight2026}.  For every finite group
$G$, scalar extension identifies
$Z(\ell G)=\ell\otimes_k Z(kG)$.  Each local factor of the finite
commutative algebra $Z(kG)$ has nilpotent radical and residue field
$k$.  After extending to $\ell$, the extended radical is a nilpotent
ideal with quotient $\ell$, so the factor remains local.
Consequently every primitive central idempotent $\bar b$ of $\ell G$
is uniquely of the form $1\otimes\bar b_0$ for a primitive central
idempotent $\bar b_0$ of $kG$.  For every $p$-subgroup $Q\leq G$,
restriction of coefficients to $C_G(Q)$ gives
\[
 \operatorname{Br}^{\ell}_Q(\bar b)
   =1\otimes\operatorname{Br}^{k}_Q(\bar b_0).
\]
Faithfulness of field extension shows that these Brauer images are
nonzero for the same $Q$.  The characterization of defect groups as
the maximal $p$-subgroups with nonzero Brauer image therefore gives the
same defect groups for the two residue blocks.

For either coefficient pair $(\Lambda,F)=(\cO,k)$ or
$(\mathcal R,\ell)$, the class sums give a $\Lambda$-basis of
$Z(\Lambda G)$ whose reductions give an $F$-basis of $Z(FG)$.
Thus $Z(\Lambda G)$ is a finite free complete commutative
$\Lambda$-algebra with residue algebra exactly $Z(FG)$.  Idempotents
in this residue algebra lift uniquely: the derivative $2x-1$ of
$x^2-x$ is a unit at an idempotent modulo the maximal ideal of
$\Lambda$, so the complete-ring Hensel argument applies, including
when $p=2$.  Primitivity is preserved, since central decompositions lift and
an idempotent reducing to zero is zero.

Now let $b$ be a block idempotent of $\mathcal R G$, and let $b_0$ be
the unique central idempotent of $\cO G$ lifting the corresponding
$\bar b_0$.  The image of $b_0$ in $\mathcal R G$ is a central
idempotent reducing to $\bar b$, so uniqueness of the central lift
makes it equal to $b$.  In particular
\begin{equation}\label{eq:complete-dvr-block-base-change}
 \mathcal R G b\cong\mathcal R\otimes_{\cO}(\cO G b_0)
\end{equation}
as $\mathcal R$-algebras, and $b_0$ has the same defect group as $b$.

Finally, scalar extension preserves the specified linear Morita
equivalences.  Indeed, for finite $\cO$-algebras $A,B$, a
$\cO$-linear Morita equivalence is represented by inverse
$\cO$-central Morita bimodules ${}_B P_A$ and ${}_A Q_B$ with
\[
 P\otimes_A Q\cong B,\qquad Q\otimes_B P\cong A.
\]
Write $A_{\mathcal R}=\mathcal R\otimes_{\cO}A$ and similarly for
$B,P,Q$.  Base change of the bimodules and their context maps gives
\[
\begin{aligned}
 P_{\mathcal R}\otimes_{A_{\mathcal R}}Q_{\mathcal R}
   &\cong\mathcal R\otimes_{\cO}(P\otimes_A Q)
    \cong B_{\mathcal R},\\
 Q_{\mathcal R}\otimes_{B_{\mathcal R}}P_{\mathcal R}
   &\cong\mathcal R\otimes_{\cO}(Q\otimes_B P)
    \cong A_{\mathcal R}.
\end{aligned}
\]
The Morita-context identities are preserved, as are finite
projectivity and the generator property.  These bimodules therefore
give an $\mathcal R$-linear Morita equivalence on finitely generated
modules.

Choose the finite list of $\cO$-block representatives of defect order
at most $M$ from Theorem~\ref{thm:main}.  The block $\cO G b_0$ in
\eqref{eq:complete-dvr-block-base-change} has that defect bound, so it
is $\cO$-linearly Morita equivalent to one representative.  The
base-changed context puts $\mathcal R G b$ in the Morita class of its
scalar extension.  These scalar extensions form a finite list over the
fixed ring $\mathcal R$, proving the corollary.
\end{proof}
